\chapter{Application}

We now demonstrate how the interventional effects developed in the previous section 
can be deployed by applying them to the example introduced in 
Section~\ref{section:vr_int_eff}.~\cite{kuhaetal2021} used data from three British
cohort studies for individuals born in 1946, 1958, and 1970 to investigate the dynamics
of intergenerational social class mobility in Britain.~\cite{kuhaetal2021}
used these cohort studies to investigate the role of level of education in mediating 
intergenerational social class mobility. They do this by using methods adapted from mediation
analysis to estimate population associations of variables. We use the same data but
additionally include the previously unused results of an intelligence test taken by 
study participants at age 10. This allows us to characterize the intergenerational 
social class mobility problem as one in which two causally ordered mediators, intelligence 
and education, mediate the relationship between a class of origin --- the parents' social 
class --- and a class of destination --- a participant's social class at mid-life. The causal 
structure we assume is encoded in Figure~\ref{fig:application_DAG}. Additionally, 
instead of associations, we estimate VanderWeele-Robinson interventional effects. 

For simplicity, we focus on men in the 1970 cohort study~\citep{cls2026} and do not consider any covariates.
For each observation, we use the following variables:

\begin{itemize}[noitemsep,topsep=0pt]
\item  The \textit{class of origin} is proxied by the participant's father's recorded social class. 
Social class is categorized using the NS-SEC classification~\citep{onsnd} for social class as determined by 
occupation, with 7 different categorizations. The categorizations are: 
\begin{enumerate}[noitemsep,topsep=0pt]
    \item Higher managers and professionals
    \item Lower managers and professionals
    \item Intermediate occupations
    \item Small employers and own account workers
    \item Lower supervisory and technical occupations
    \item Semiroutine occupations
    \item Routine occupations 
\end{enumerate}
This is generally regarded as an unordered categorical 
variable~\citep{kuhaetal2021}. ~\cite{kuhaetal2021}, 


\item  The \textit{class of destination} is equal to the participant's recorded social class at middle 
age (age 38). 
\item Intelligence is measured by a teacher-adapted version of the British Ability Scales~\citep{hill2005} 
that was administered in schools. Four subscales --- word definitions, word similarities, 
recall of digits, and matrices --- were used, and a principal components analysis was carried out on the 
resulting data to derive a general cognitive ability factor ($g$)~\citep{galeetal2007}. This score was 
then standardized into a z-score to measure ability relative to the rest of the cohort. 
\item  Education is coded in terms of relative education as defined by~\cite{bukodigoldthorpe2016}. This is
done by recording an individual's highest level of education attained at age 37 and then grouping the 
individual's level of education relative to their cohort. Thus, this is an ordered categorical variable,
which we treat as a continuous variable for simplicity. 
\end{itemize}

There are several important caveats for our analysis. These cohort studies are observational data, 
and other than the treatment, mediator, and outcome variables, we have few other possible covariates. 
This means that the assumptions required to identify even VanderWeele-Robinson interventional effects 
are not fulfilled, as there are many confounders we would need to control to even plausibly argue
that these assumptions are satisfied. Thus, the numbers offered in this analysis should be 
interpreted with caution and cannot readily admit a causal interpretation. The main point of our 
analysis is to demonstrate how such a multiple mediator mediation analysis problem can be reasoned 
about using interventional effects. Additionally, we include the other interventional effect measures
discussed in this dissertation as points of comparison. As will be discussed below, these estimates 
are entirely illustrative as the treatment (social class) is non-manipulable. However, it is instructive 
to compare the numerical estimates from these methods to the ``main'' VanderWeele-Robinson estimates 
we present. Finally, only about 45\% of all observations for men in 1970 have complete data for 
origin and destination class, IQ, and education; unlike~\cite{kuhaetal2021}, who addresses this issue
via multiple imputation, we conduct a complete case analysis where we drop any observations with missing 
data. Overall, this leaves us with 2,730 observations in our dataset. 

We use Monte Carlo estimation to estimate interventional effects. To do this, we fit successive linear models
for intelligence, education, and child's social class as functions of the variables that causally
precede them in R~\citep{areletal2024}. We then sample many independent draws from the fitted conditional 
distributions and apply
a sequential sampling strategy as needed to estimate the needed g-formula equations for interventional effects
(e.g. Equations~\ref{eq:RIE_one_mediator_g_comp},~\ref{eq:RIE_two_mediator_g_comp},~\ref{eq:g_formula_vr_1m}). 
Finally, we use the nonparametric bootstrap for inference and construct percentile bootstrap 
confidence intervals. A link to the R code used to both estimate the effects shown in this chapter and generate
the associated tables can be found
\href{https://github.com/hannhat/personal/tree/main/research/MSc_Thesis-Mediation_Analysis/code/analysis}{here}. 

Such a strategy has been also implemented 
in~\cite{imaietal2010a} to estimate causal mediation effects, and similar strategies are described 
in~\cite{vansteelandtdaniel2017} and~\cite{kuhaetal2021}. In general, this method allows for fast and 
flexible specification of unique models for the conditional distributions of the mediators and outcome 
without the need to re-derive expressions for each path-specific interventional effect in terms of the 
regression coefficients~\citep{vansteelandtdaniel2017}. 

Here, we specify linear
models for the mediators and the outcome. The outcome of interest for our main specification is the 
probability of the participant's social class at middle age being Class 1. We observe a large disparity in 
the probability of the participant's social class being Class 1 if their 
father was in Class 1 compared to if their father was in Class 7. We are therefore interested in what 
effect various interventions to equalize the distributions of IQ and education for participants with Class 7
fathers to those with Class 1 fathers would have on their probability of ending up in Class 1. Throughout 
this analysis, following similar analyses in~\cite{jacksonvanderweele2018}, our binary outcomes are expressed 
on the risk difference scale. Thus, the observed disparity of $0.301$ between children of Class 1 compared
to Class 7 in terms of their probability of ending up in Class 1 reflects the fact that children of fathers in 
Class 7 have a baseline probability (12\%) of ending in Class 1, while children of fathers in Class 1 have 
probability to end in Class 1 equal to the baseline probability plus 30\% (42\%). 

In our main specification (Table~\ref{tab:app_main_spec}), we specify successive, fully additive linear models for the mediators and the 
outcome, with IQ a function of father's class, education a function of father's class and IQ, and participant's 
class a function of father's class, education, and IQ. We find that the interventions to equalize the marginal 
distribution of IQ only and the marginal 
distribution of education only result in a 25\% and 23\% reductions, respectively, in the difference in probability 
of ending up in Class 1; the intervention to equalize the distribution of education within levels of IQ results 
in a 15\% reduction; and the intervention to equalize the joint distribtion of IQ and education results in a 
42\% reduction. All estimates are statistically significant relative to a null hypothesis of no effect at the 
95\% confidence interval. 

\begin{table}[ht]
\centering
\sbox0{%
\begin{tabular}{clcc}
\hline
\addlinespace[6pt]
& & Estimate & \shortstack{95\% Confidence \\ Interval} \\
\addlinespace[3pt]
\noalign{\hrule height 1pt}
\addlinespace[3pt]
1 & \makecell[l]{Total disparity} & 0.301 & [0.242, 0.359] \\
\addlinespace[2pt]
2 & \makecell[l]{Effect of equalizing IQ} & 0.076 & [0.050, 0.099] \\
\addlinespace[2pt]
3 & \makecell[l]{Effect of equalizing education} & 0.068 & [0.045, 0.093] \\
\addlinespace[2pt]
4 & \makecell[l]{Effect of equalizing education,\\ conditional on IQ} & 0.046 & [0.026, 0.068] \\
\addlinespace[2pt]
5 & \makecell[l]{Effect of jointly equalizing\\ IQ and education} & 0.126 & [0.097, 0.154] \\
\addlinespace[3pt]
\noalign{\hrule height 1pt}
\end{tabular}%
}%
\captionsetup{width=\wd0}
\caption{The effect of VanderWeele-Robinson interventions to equalize mediators between
Class 7 and Class 1 on Class 7 children's probabilities of ending in Class 1  --- fully additive linear models}
\label{tab:app_main_spec}
\usebox0 \\[4pt]
\parbox{\wd0}{\footnotesize
\emergencystretch=3em
\setlength{\parindent}{1.5em}
\textsuperscript{1} Estimates were constructed using 1,000,000 Monte Carlo samples and 500 bootstrap 
replications.\par}
\end{table}

We also run an alternative specification (Table~\ref{tab:app_interaction_spec}) where we specify successive linear models for IQ, education,
and participant's social class but now include all possible mediator-mediator and mediator-outcome interactions 
in our models. We find broadly similar effect estimates even when including these interactions for the 
interventions on single mediators, but we do find a substantial increase in the effect of the joint 
intervention to equalize IQ and education: this intervention now leads to a 50\% reduction in the disparity 
in the probability of reaching Class 1 between Class 7 and Class 1 children. 

\begin{table}[ht]
\centering
\sbox0{%
\begin{tabular}{clcc}
\hline
\addlinespace[6pt]
& & Estimate & \shortstack{95\% Confidence \\ Interval} \\
\addlinespace[3pt]
\noalign{\hrule height 1pt}
\addlinespace[3pt]
1 & \makecell[l]{Total disparity} & 0.301 & [0.241, 0.359] \\
\addlinespace[2pt]
2 & \makecell[l]{Effect of equalizing IQ} & 0.078 & [0.048, 0.106] \\
\addlinespace[2pt]
3 & \makecell[l]{Effect of equalizing education} & 0.064 & [0.033, 0.097] \\
\addlinespace[2pt]
4 & \makecell[l]{Effect of equalizing education,\\ conditional on IQ} & 0.045 & [0.021, 0.068] \\
\addlinespace[2pt]
5 & \makecell[l]{Effect of jointly equalizing\\ IQ and education} & 0.152 & [0.115, 0.188] \\
\addlinespace[3pt]
\noalign{\hrule height 1pt}
\end{tabular}%
}%
\captionsetup{width=\wd0}
\caption{The effect of VanderWeele-Robinson interventions to equalize mediators between
Class 7 and Class 1 on Class 7 children's probabilities of ending in Class 1  --- fully interacted linear models}
\label{tab:app_interaction_spec}
\usebox0 \\[4pt]
\parbox{\wd0}{\footnotesize
\emergencystretch=3em
\setlength{\parindent}{1.5em}
\textsuperscript{1} Estimates were constructed using 1,000,000 Monte Carlo samples and 500 bootstrap 
replications.\par}
\end{table}

Finally, we conduct an analysis where we now consider a hierarchical ordering of social classes. 
Following~\cite{kuhaetal2021}, we take Classes 3, 4, and 5 to be on the same level, so classes $c = {1,2,3,4,5,6,7}$
now map onto hierarchical positions $h_{c} = {1,2,3,3,3,4,5}$. We then consider child's social class as a continuous
outcome variable instead of a binary variable. Thus, we now analyze how the disparity in the average hierarchical
social class position changes for children of Class 7 relative to children of Class 1 under various interventions
to equalize IQ and education among these social classes (Table~\ref{tab:app_hierarchical_spec}). We specify successive,
fully additive linear models with no interactions for IQ, education, and child's social class. We find 
% insert interpretation of coefficients here

\begin{table}[ht]
\centering
\sbox0{%
\begin{tabular}{clcc}
\hline
\addlinespace[6pt]
& & Estimate & \shortstack{95\% Confidence \\ Interval} \\
\addlinespace[3pt]
\noalign{\hrule height 1pt}
\addlinespace[3pt]
1 & \makecell[l]{Total disparity} & -1.126 & [-1.306, -0.977] \\
\addlinespace[2pt]
2 & \makecell[l]{Effect of equalizing IQ} & -0.344 & [-0.422, -0.262] \\
\addlinespace[2pt]
3 & \makecell[l]{Effect of equalizing education} & -0.378 & [-0.457, -0.295] \\
\addlinespace[2pt]
4 & \makecell[l]{Effect of equalizing education,\\ conditional on IQ} & -0.236 & [-0.312, -0.164] \\
\addlinespace[2pt]
5 & \makecell[l]{Effect of jointly equalizing\\ IQ and education} & -0.565 & [-0.659, -0.474] \\
\addlinespace[3pt]
\noalign{\hrule height 1pt}
\end{tabular}%
}%
\captionsetup{width=\wd0}
\caption{The effect of VanderWeele-Robinson interventions to
equalize mediators between
Class 7 and Class 1 on Class 7 children's average hierarchical class position}
\label{tab:app_hierarchical_spec}
\usebox0 \\[4pt]
\parbox{\wd0}{\footnotesize
\emergencystretch=3em
\setlength{\parindent}{1.5em}
\textsuperscript{1} Estimates were constructed using 1,000,000 Monte Carlo samples and 500 bootstrap 
replications.\par}
\end{table}

Similar to the results of~\cite{kuhaetal2021}, the results of our binary outcome specifications, 
if causally identified, would provide suggestive 
evidence that an intervention to make educational distributions between social classes more equal might not 
make a large difference in reducing the observed discrepancies in intergenerational social class mobility. Any
similar hypothetical intervention on setting the marginal distribution of intelligence would have about the 
same effect as intervening to equalize the distribution of education. Finally, an intervention to equalize the 
joint distribution of IQ and education would lead to a 40-50\% reduction in the social class mobility disparity, 
which would be a more significant reduction if such an intervention were feasible and ethical. 

As we have argued in the previous sections, given the purely prescriptive interpretation of these effects, the 
usefulness of these estimates is guided by the feasibility of the posited hypothetical interventions for the 
substantive problem of interest. In this example, hypothetical interventions to equalize levels of education could be 
feasibly approximated. For example, if we believe a significant source of the disparity in levels of education is due 
to more favored social classes (e.g. Class 1) being able to spend more money to send their children to better schools, 
a proposal to require attendance at state schools that can be tested into via standardized testing might approximate an 
intervention to equalize education conditional on intelligence. A more involved quota system whereby an equal proportion 
of students from each social class are alotted spots in schools might correspond to an intervention to equalize the 
marginal distributions of education. Interventions to affect the distribution of intelligence, meanwhile, are both more 
difficult to conceive of and ethically less palatable. 

Thus, when weighing the practical implications of these effects, we might be more interested in how effective 
the interventions to equalize only education might be. Suppose we want to get a more complete picture of how these effects
would affect the probability of reaching Class 1 for all social classes under a variety of counterfactual interventions to 
set their education distribution to that of other social classes. 

\begin{table}[ht]
\centering
\sbox0{%
\begin{tabular}{ccrrrrrrr}
\hline
\addlinespace[3pt]
& & \multicolumn{7}{c}{Class} \\
\addlinespace[3pt]
& & \multicolumn{1}{c}{(1)} & \multicolumn{1}{c}{(2)} & \multicolumn{1}{c}{(3)} & \multicolumn{1}{c}{(4)} & \multicolumn{1}{c}{(5)} & \multicolumn{1}{c}{(6)} & \multicolumn{1}{c}{(7)} \\
\addlinespace[3pt]
\noalign{\hrule height 1pt}
\addlinespace[3pt]
\multirow{7}{*}{\shortstack{Class}}
& (1) & \multicolumn{1}{c}{0.000} & 0.114 & 0.153 & 0.238 & 0.191 & 0.263 & 0.301 \\
& (2) & & \multicolumn{1}{c}{0.000} & 0.039 & 0.124 & 0.077 & 0.149 & 0.187 \\
& (3) & & & \multicolumn{1}{c}{0.000} & 0.086 & 0.038 & 0.110 & 0.149 \\
& (4) & & & & \multicolumn{1}{c}{0.000} & -0.048 & 0.025 & 0.063 \\
& (5) & & & & & \multicolumn{1}{c}{0.000} & 0.072 & 0.111 \\
& (6) & & & & & & \multicolumn{1}{c}{0.000} & 0.038 \\
& (7) & & & & & & & \multicolumn{1}{c}{0.000} \\
\addlinespace[3pt]
\noalign{\hrule height 1pt}
\end{tabular}%
}%
\captionsetup{width=\wd0}
\caption{Differences in the probability of ending up in Class 1 given parents of different classes}
\label{tab:app_disparity_matrix}
\usebox0 \\[4pt]
\parbox{\wd0}{\footnotesize
\emergencystretch=3em
\setlength{\parindent}{1.5em}
\textsuperscript{1} Raw disparity = row class - column class.\par}
\end{table}

To help visualize what the effects of such interventions might look like, we present a matrix of changes 
in the probability of ending up in Class 1 for interventions for a reference class to equalize both marginal 
(Table~\ref{tab:app_marginal_education_matrix}) and conditional (Table~\ref{tab:app_conditional_education_matrix}) distributions of 
education to those of another class. Table~\ref{tab:app_disparity_matrix} displays the raw
disparities in probability of ending in Class 1 for all combinations of classes. Compared to the previous tables,
which gave the effects of equalizing the distribution of education for Class 7 to that of Class 1 only, 
these tables give a more complete account of the effects of all possible equalizing interventions. 

As we can see from Tables~\ref{tab:app_marginal_education_matrix} and~\ref{tab:app_conditional_education_matrix}, the more aggressive 
intervention to equalize education marginally as opposed to
within levels of IQ is a fair amount more impactful in increasing (or decreasing, in the case of ``higher'' social classes
being set to the education distribution of ``lower'' ones) the probability of ending up in Class 1. An additional point of caution
for interpreting these numbers, even in a hypothetical scenario in which we had properly controlled for all other unmeasured confounders,
the effect of equalizing the marginal distribution of education would still not be identified because IQ would be an unmeasured confounder
for this effect (see Section~\ref{subsection:vr_effects_2m}). Thus, actually directly comparing these two potential interventions would 
require an experimental setup where children from different origin classes were randomly assigned levels of education from another class. 
If such an experimental setup were feasible, or if given a method to bound the effect of unmeasured confounding from IQ, these numbers could be 
used as a data point to determine the benefits of proposed policies compared with the cost of implementing them and other 
practical and ethical considerations. 

\begin{table}[ht]
\centering
\sbox0{%
\begin{tabular}{ccrrrrrrr}
\hline
\addlinespace[3pt]
& & \multicolumn{7}{c}{Class to draw level of education from} \\
\addlinespace[3pt]
& & \multicolumn{1}{c}{(1)} & \multicolumn{1}{c}{(2)} & \multicolumn{1}{c}{(3)} & \multicolumn{1}{c}{(4)} & \multicolumn{1}{c}{(5)} & \multicolumn{1}{c}{(6)} & \multicolumn{1}{c}{(7)} \\
\addlinespace[3pt]
\noalign{\hrule height 1pt}
\addlinespace[3pt]
\multirow{7}{*}{\shortstack{Reference \\ class}}
& (1) & -0.049* & -0.068* & -0.087* & -0.105* & -0.124* & -0.141* & -0.158* \\
& (2) & 0.029\phantom{*} & 0.011\phantom{*} & -0.007\phantom{*} & -0.025\phantom{*} & -0.042* & -0.058* & -0.074* \\
& (3) & 0.027\phantom{*} & 0.009\phantom{*} & -0.008\phantom{*} & -0.024\phantom{*} & -0.039\phantom{*} & -0.054* & -0.068* \\
& (4) & 0.073* & 0.057* & 0.042* & 0.027\phantom{*} & 0.013\phantom{*} & 0.000\phantom{*} & -0.013\phantom{*} \\
& (5) & 0.009\phantom{*} & -0.006\phantom{*} & -0.021\phantom{*} & -0.035\phantom{*} & -0.049* & -0.061* & -0.073* \\
& (6) & 0.051* & 0.037* & 0.024\phantom{*} & 0.011\phantom{*} & -0.001\phantom{*} & -0.012\phantom{*} & -0.023\phantom{*} \\
& (7) & 0.068* & 0.055* & 0.043* & 0.031* & 0.020* & 0.010\phantom{*} & 0.000\phantom{*} \\
\addlinespace[3pt]
\noalign{\hrule height 1pt}
\end{tabular}%
}%
\captionsetup{width=\wd0}
\caption{Change in probability of ending in Class 1 between
two classes induced by setting the marginal distribution of education
for one class (row) to that of the other (column)}
\label{tab:app_marginal_education_matrix}
\usebox0 \\[4pt]
\parbox{\wd0}{\footnotesize}
\emergencystretch=3em
\setlength{\parindent}{1.5em}
\textsuperscript{1} Each number represents the estimated change in probability of ending in Class 1 as a result of 
drawing education for the reference class from the class indicated in the column.\par
\textsuperscript{2} Estimates of the interventional effect were constructed using 1,000,000 Monte Carlo samples and 100 
bootstrap replications. The presence of a significance star indicates that the 95\% bootstrap confidence interval for the 
interventional effect did not contain zero. \par
\end{table}

\begin{table}[ht]
\centering
\sbox0{%
\begin{tabular}{ccrrrrrrr}
\hline
\addlinespace[3pt]
& & \multicolumn{7}{c}{Class to draw level of education from} \\
\addlinespace[3pt]
& & \multicolumn{1}{c}{(1)} & \multicolumn{1}{c}{(2)} & \multicolumn{1}{c}{(3)} & \multicolumn{1}{c}{(4)} & \multicolumn{1}{c}{(5)} & \multicolumn{1}{c}{(6)} & \multicolumn{1}{c}{(7)} \\
\addlinespace[3pt]
\noalign{\hrule height 1pt}
\addlinespace[3pt]
\multirow{7}{*}{\shortstack{Reference \\ class}}
& (1) & -0.041* & -0.052* & -0.063* & -0.074* & -0.085* & -0.095* & -0.105* \\
& (2) & 0.029\phantom{*} & 0.019\phantom{*} & 0.009\phantom{*} & -0.002\phantom{*} & -0.012\phantom{*} & -0.021\phantom{*} & -0.031\phantom{*} \\
& (3) & 0.015\phantom{*} & 0.006\phantom{*} & -0.004\phantom{*} & -0.013\phantom{*} & -0.022\phantom{*} & -0.031\phantom{*} & -0.039\phantom{*} \\
& (4) & 0.054* & 0.045* & 0.036* & 0.029\phantom{*} & 0.021\phantom{*} & 0.013\phantom{*} & 0.005\phantom{*} \\
& (5) & -0.008\phantom{*} & -0.017\phantom{*} & -0.025\phantom{*} & -0.033\phantom{*} & -0.041* & -0.048* & -0.055* \\
& (6) & 0.030* & 0.023\phantom{*} & 0.016\phantom{*} & 0.009\phantom{*} & 0.002\phantom{*} & -0.005\phantom{*} & -0.011\phantom{*} \\
& (7) & 0.046* & 0.039* & 0.033* & 0.026* & 0.020* & 0.014\phantom{*} & 0.009\phantom{*} \\
\addlinespace[3pt]
\noalign{\hrule height 1pt}
\end{tabular}%
}%
\captionsetup{width=\wd0}
\caption{Change in probability of ending in Class 1 between
two classes induced by setting the distribution of education, conditional on IQ,
for one class (row) to that of the other (column)}
\label{tab:app_conditional_education_matrix}
\usebox0 \\[4pt]
\parbox{\wd0}{\footnotesize}
\emergencystretch=3em
\setlength{\parindent}{1.5em}
\textsuperscript{1} Each number represents the estimated change in probability of ending in Class 1 as a result of 
drawing education for the reference class from the class indicated in the column.\par
\textsuperscript{2} Estimates of the interventional effect were constructed using 1,000,000 Monte Carlo samples and 100 
bootstrap replications. The presence of a significance star indicates that the 95\% bootstrap confidence interval for the 
interventional effect did not contain zero. \par
\end{table}